\textbf{Phần 1. Trắc nghiệm lựa chọn}
\setcounterpageref{ex}{0}

\begin{ex} %[5D2V13]
Một số thập phân có hai chữ số ở phần thập phân. Sau khi làm tròn đến hàng phần mười thì được số $8{,}6$. Số thập phân đó lớn nhất có thể là số nào sau đây?
\choice
{$8{,}69$}
{\True $8{,}64$}
{$8{,}65$}
{$8{,}59$}
\loigiai{
Để làm tròn đến hàng phần mười được $8{,}6$:
\begin{itemize}
    \item Nếu làm tròn xuống thì phần nguyên là $8$, hàng phần mười giữ nguyên là $6$, và chữ số hàng phần trăm lớn nhất phải nhỏ hơn $5$, tức là $4$. Số lớn nhất là $8{,}64$.
    \item Nếu làm tròn lên thì hàng phần mười là $5$ và hàng phần trăm là $5, 6, 7, 8, 9$, lớn nhất là $8{,}59 < 8{,}64$.
\end{itemize}
Vậy số lớn nhất thỏa mãn là $8{,}64$.\\
\textbf{Đáp án của câu hỏi là B}
}
\end{ex}

\begin{ex} %[5D2V13]
Một số thập phân có hai chữ số ở phần thập phân. Sau khi làm tròn đến hàng phần mười thì được số $5{,}3$. Số thập phân đó bé nhất có thể là số nào sau đây?
\choice
{$5{,}31$}
{$5{,}30$}
{\True $5{,}25$}
{$5{,}20$}
\loigiai{
Để làm tròn lên được $5{,}3$ thì hàng phần mười phải là $2$ và chữ số hàng phần trăm nhỏ nhất phải là $5$.\\
Do đó số bé nhất thỏa mãn là $5{,}25$.\\
\textbf{Đáp án của câu hỏi là C}
}
\end{ex}

\begin{ex} %[5D2V13]
Tìm chữ số $x$ để khi làm tròn số thập phân $15{,}x6$ đến số tự nhiên gần nhất ta được số $16$:
\choice
{$x \in \{0, 1, 2, 3, 4\}$}
{\True $x \in \{5, 6, 7, 8, 9\}$}
{$x = 5$}
{$x = 6$}
\loigiai{
Để làm tròn lên thành $16$, chữ số ở hàng phần mười phải lớn hơn hoặc bằng $5$.\\
Do đó $x \in \{5, 6, 7, 8, 9\}$.\\
\textbf{Đáp án của câu hỏi là B}
}
\end{ex}

\begin{ex} %[5D2V13]
Điểm trung bình ba môn Toán, Tiếng Việt, Tiếng Anh của bạn Hà là $8{,}666\dots$ Nếu làm tròn điểm trung bình này đến hàng phần mười thì điểm của bạn Hà là
\choice
{$8{,}6$}
{\True $8{,}7$}
{$8{,}67$}
{$9{,}0$}
\loigiai{
Xét hàng phần trăm có chữ số $6 \ge 5$, do đó làm tròn lên hàng phần mười ta được $8{,}7$.\\
\textbf{Đáp án của câu hỏi là B}
}
\end{ex}

\begin{ex} %[5D2V13]
Khối lượng của một kiện hàng xuất khẩu là $48{,}375\text{ kg}$. Nhân viên kiểm định làm tròn khối lượng này đến hàng phần trăm. Khối lượng kiện hàng sau khi làm tròn là
\choice
{$48{,}37\text{ kg}$}
{\True $48{,}38\text{ kg}$}
{$48{,}4\text{ kg}$}
{$48\text{ kg}$}
\loigiai{
Xét chữ số hàng phần nghìn là $5 \ge 5$, nên làm tròn lên hàng phần trăm ta cộng thêm $1$ vào chữ số hàng phần trăm: $7 + 1 = 8$.\\
Kết quả sau khi làm tròn là $48{,}38\text{ kg}$.\\
\textbf{Đáp án của câu hỏi là B}
}
\end{ex}

\begin{ex} %[5D2V13]
Có bao nhiêu số thập phân có hai chữ số ở phần thập phân mà khi làm tròn đến số tự nhiên gần nhất thì được số $7$?
\choice
{$99$}
{\True $100$}
{$101$}
{$90$}
\loigiai{
Các số làm tròn thành $7$ gồm:
\begin{itemize}
    \item Các số làm tròn lên: từ $6{,}50$ đến $6{,}99$ (gồm $50$ số).
    \item Các số làm tròn xuống: từ $7{,}00$ đến $7{,}49$ (gồm $50$ số).
\end{itemize}
Tổng số các số thỏa mãn là:
\[
50 + 50 = 100\text{ (số)}.
\]
\textbf{Đáp án của câu hỏi là B}
}
\end{ex}

\begin{ex} %[5D2V13]
Nhiệt độ ngoài trời lúc $12$ giờ trưa đo được là $36{,}84^\circ\text{C}$. Nếu làm tròn nhiệt độ này đến số tự nhiên gần nhất thì nhiệt độ khoảng
\choice
{$36^\circ\text{C}$}
{\True $37^\circ\text{C}$}
{$36{,}8^\circ\text{C}$}
{$36{,}5^\circ\text{C}$}
\loigiai{
Chữ số hàng phần mười là $8 \ge 5$, nên ta làm tròn lên thành $37^\circ\text{C}$.\\
\textbf{Đáp án của câu hỏi là B}
}
\end{ex}

\begin{ex} %[5D2V13]
Một người đi xe đạp đi được quãng đường $12{,}64\text{ km}$ trong giờ thứ nhất và $11{,}85\text{ km}$ trong giờ thứ hai. Làm tròn tổng quãng đường người đó đã đi trong cả hai giờ đến hàng phần mười ta được
\choice
{$24{,}4\text{ km}$}
{\True $24{,}5\text{ km}$}
{$24{,}49\text{ km}$}
{$25\text{ km}$}
\loigiai{
Tổng quãng đường thực tế người đó đi được là:
\[
12{,}64 + 11{,}85 = 24{,}49\text{ (km)}.
\]
Xét số $24{,}49$: chữ số hàng phần trăm là $9 \ge 5$, do đó làm tròn lên hàng phần mười ta được $24{,}5\text{ km}$.\\
\textbf{Đáp án của câu hỏi là B}
}
\end{ex}

\begin{ex} %[5D2V13]
Tìm chữ số $y$ để số thập phân $9{,}8y4$ khi làm tròn đến hàng phần mười thì được số $9{,}8$:
\choice
{$y \in \{5, 6, 7, 8, 9\}$}
{\True $y \in \{0, 1, 2, 3, 4\}$}
{$y = 4$}
{$y = 0$}
\loigiai{
Để làm tròn xuống thành $9{,}8$ thì chữ số hàng phần trăm $y$ phải nhỏ hơn $5$.\\
Do đó $y \in \{0, 1, 2, 3, 4\}$.\\
\textbf{Đáp án của câu hỏi là B}
}
\end{ex}

\begin{ex} %[5D2V13]
Một thửa ruộng thu hoạch được $2{,}486\text{ tấn}$ thóc. Sau khi làm tròn số thóc đến hàng phần mười, bác nông dân nói: ``Vụ này tôi thu được khoảng $2{,}5\text{ tấn}$ thóc''. Hỏi bác nông dân đã làm tròn lên bao nhiêu phần nghìn của một tấn thóc?
\choice
{$0{,}014\text{ tấn}$}
{\True $\dfrac{14}{1\,000}\text{ tấn}$}
{$\dfrac{14}{100}\text{ tấn}$}
{$\dfrac{4}{1\,000}\text{ tấn}$}
\loigiai{
Độ chênh lệch giữa số sau khi làm tròn và số thực tế là:
\[
2{,}5 - 2{,}486 = 0{,}014\text{ (tấn)} = \dfrac{14}{1\,000}\text{ tấn}.
\]
\textbf{Đáp án của câu hỏi là B}
}
\end{ex}

\begin{ex} %[5D2V13]
Số thập phân $19{,}95$ khi làm tròn đến hàng phần mười sẽ được số nào sau đây?
\choice
{$19{,}9$}
{\True $20{,}0$}
{$20$}
{$19{,}10$}
\loigiai{
Chữ số hàng phần trăm là $5 \ge 5$, nên làm tròn lên: chữ số hàng phần mười $9 + 1 = 10$, viết $0$ nhớ $1$ sang phần nguyên được $20{,}0$.\\
\textbf{Đáp án của câu hỏi là B}
}
\end{ex}

\begin{ex} %[5D2V13]
Một chiếc xe chở $4$ thùng hàng có khối lượng lần lượt là: $12{,}3\text{ kg}$; $15{,}8\text{ kg}$; $18{,}4\text{ kg}$; $21{,}6\text{ kg}$. Nếu làm tròn khối lượng mỗi thùng đến số tự nhiên gần nhất rồi tính tổng thì khối lượng ước tính là
\choice
{$67\text{ kg}$}
{\True $68\text{ kg}$}
{$69\text{ kg}$}
{$68{,}1\text{ kg}$}
\loigiai{
Làm tròn từng thùng đến số tự nhiên gần nhất:
\begin{itemize}
    \item $12{,}3 \approx 12\text{ kg}$.
    \item $15{,}8 \approx 16\text{ kg}$.
    \item $18{,}4 \approx 18\text{ kg}$.
    \item $21{,}6 \approx 22\text{ kg}$.
\end{itemize}
Tổng ước tính là: $12 + 16 + 18 + 22 = 68\text{ (kg)}$.\\
\textbf{Đáp án của câu hỏi là B}
}
\end{ex}

\textbf{Phần 2. Trắc nghiệm đúng sai}
\setcounterpageref{ex}{0}

\begin{ex} %[5D2V13]
Cho số thập phân $P = 47{,}648$. Xét tính đúng sai của các khẳng định sau:
\choiceTF
{\True Làm tròn số $P$ đến số tự nhiên gần nhất ta được $48$.}
{\True Làm tròn số $P$ đến hàng phần mười ta được $47{,}6$.}
{Làm tròn số $P$ đến hàng phần trăm ta được $47{,}64$.}
{\True Hiệu giữa số $P$ và kết quả làm tròn đến hàng phần mười là $0{,}048$.}
\loigiai{
\begin{itemize}
    \item Chữ số hàng phần mười là $6 \ge 5$, làm tròn đến số tự nhiên được $48$. Khẳng định a \textbf{Đúng}.
    \item Chữ số hàng phần trăm là $4 < 5$, làm tròn đến hàng phần mười được $47{,}6$. Khẳng định b \textbf{Đúng}.
    \item Chữ số hàng phần nghìn là $8 \ge 5$, làm tròn đến hàng phần trăm phải được $47{,}65 
e 47{,}64$. Khẳng định c \textbf{Sai}.
    \item $47{,}648 - 47{,}6 = 0{,}048$. Khẳng định d \textbf{Đúng}.
\end{itemize}
\begin{itemize}
    \item a) \textbf{Đúng}.
    \item b) \textbf{Đúng}.
    \item c) \textbf{Sai}: Kết quả là $47{,}65$.
    \item d) \textbf{Đúng}.
\end{itemize}
}
\end{ex}

\begin{ex} %[5D2V13]
Số lượng sữa trong ba hộp lần lượt là: $0{,}95\text{ l}$; $1{,}24\text{ l}$ và $0{,}82\text{ l}$. Xét tính đúng sai của các khẳng định sau:
\choiceTF
{\True Làm tròn số lượng sữa của hộp thứ nhất đến số tự nhiên gần nhất ta được $1\text{ l}$.}
{Làm tròn số lượng sữa của hộp thứ hai đến hàng phần mười ta được $1{,}3\text{ l}$.}
{\True Tổng lượng sữa thực tế trong cả ba hộp là $3{,}01\text{ l}$.}
{\True Làm tròn tổng lượng sữa thực tế đến hàng phần mười ta được $3{,}0\text{ l}$.}
\loigiai{
\begin{itemize}
    \item $0{,}95$ có hàng phần mười là $9$, làm tròn lên được $1\text{ l}$. Khẳng định a \textbf{Đúng}.
    \item $1{,}24$ có hàng phần trăm là $4 < 5$, làm tròn đến hàng phần mười được $1{,}2\text{ l} 
e 1{,}3\text{ l}$. Khẳng định b \textbf{Sai}.
    \item Tổng thực tế: $0{,}95 + 1{,}24 + 0{,}82 = 3{,}01\text{ (l)}$. Khẳng định c \textbf{Đúng}.
    \item $3{,}01$ có hàng phần trăm là $1 < 5$, làm tròn đến hàng phần mười được $3{,}0\text{ l}$. Khẳng định d \textbf{Đúng}.
\end{itemize}
\begin{itemize}
    \item a) \textbf{Đúng}.
    \item b) \textbf{Sai}: Làm tròn thành $1{,}2\text{ l}$.
    \item c) \textbf{Đúng}.
    \item d) \textbf{Đúng}.
\end{itemize}
}
\end{ex}

\textbf{Phần 3. Trả lời ngắn}
\setcounterpageref{ex}{0}

\begin{ex} %[5D2V13]
Một số thập phân có hai chữ số ở phần thập phân. Sau khi làm tròn đến hàng phần mười thì được số $9{,}4$. Số đó lớn nhất có thể là bao nhiêu? (Nhập dưới dạng số thập phân, ví dụ $9.44$).
\shortans{9.44}
\loigiai{
\textbf{Đáp số:} $9.44$\\
Để làm tròn đến hàng phần mười được $9{,}4$ mà là số lớn nhất:
Phần nguyên là $9$, hàng phần mười là $4$, chữ số hàng phần trăm lớn nhất phải nhỏ hơn $5$, tức là $4$.
Vậy số đó là $9{,}44$.
}
\end{ex}

\begin{ex} %[5D2V13]
Một số thập phân có hai chữ số ở phần thập phân. Sau khi làm tròn đến hàng phần mười thì được số $7{,}5$. Số đó bé nhất có thể là bao nhiêu? (Nhập dưới dạng số thập phân, ví dụ $7.45$).
\shortans{7.45}
\loigiai{
\textbf{Đáp số:} $7.45$\\
Để làm tròn lên hàng phần mười được $7{,}5$ mà là số bé nhất:
Phần nguyên là $7$, hàng phần mười là $4$, chữ số hàng phần trăm nhỏ nhất phải là $5$.
Vậy số đó là $7{,}45$.
}
\end{ex}

\begin{ex} %[5D2V13]
Làm tròn số thập phân $24{,}875$ đến hàng phần mười ta được số có dạng $24{,}a$. Hãy tìm giá trị của $a$.
\shortans{9}
\loigiai{
\textbf{Đáp số:} $9$\\
Chữ số hàng phần trăm là $7 \ge 5$, do đó khi làm tròn đến hàng phần mười ta cộng thêm $1$ vào chữ số hàng phần mười: $8 + 1 = 9$.
Số nhận được là $24{,}9$, tức $a = 9$.
}
\end{ex}

\begin{ex} %[5D2V13]
Làm tròn số thập phân $6{,}482$ đến hàng phần trăm ta được số có dạng $6{,}4b$. Hãy tìm giá trị của $b$.
\shortans{8}
\loigiai{
\textbf{Đáp số:} $8$\\
Chữ số hàng phần nghìn là $2 < 5$, do đó khi làm tròn đến hàng phần trăm ta giữ nguyên chữ số hàng phần trăm là $8$.
Vậy $b = 8$.
}
\end{ex}

\begin{ex} %[5D2V13]
Có tất cả bao nhiêu số tự nhiên $x$ thỏa mãn điều kiện: khi làm tròn số thập phân $x{,}48$ đến số tự nhiên gần nhất thì được một số thuộc tập hợp $\{10, 11, 12, 13\}$?
\shortans{4}
\loigiai{
\textbf{Đáp số:} $4$\\
Vì số có dạng $x{,}48$ có chữ số hàng phần mười là $4 < 5$, nên khi làm tròn đến số tự nhiên gần nhất ta luôn làm tròn xuống bằng chính $x$.
Do đó để kết quả làm tròn thuộc $\{10, 11, 12, 13\}$ thì $x \in \{10, 11, 12, 13\}$.
Vậy có $4$ số tự nhiên $x$ thỏa mãn.
}
\end{ex}

\begin{ex} %[5D2V13]
Tính tổng hai số $15{,}74$ và $24{,}38$, sau đó làm tròn kết quả đến hàng phần mười ta được số $40{,}c$. Hãy tìm giá trị của chữ số $c$.
\shortans{1}
\loigiai{
\textbf{Đáp số:} $1$\\
Tổng hai số là:
\[
15{,}74 + 24{,}38 = 40{,}12.
\]
Làm tròn số $40{,}12$ đến hàng phần mười: chữ số hàng phần trăm là $2 < 5$, nên giữ nguyên chữ số hàng phần mười là $1$.
Kết quả là $40{,}1$, vậy $c = 1$.
}
\end{ex}
